\section*{Chapter \ref{ch:b1:e-ph-diagrams} --- E--pH Diagrams}

\begin{solution}{exo:b1:e-ph-diagrams:1}
Manganese: \ce{Mn} (0) $<$ \ce{Mn^2+}, \ce{Mn(OH)2} ($+$II) $<$ \ce{MnO2}
($+$IV) $<$ \ce{MnO4-} ($+$VII). Chlorine: \ce{Cl-} ($-$I) $<$ \ce{Cl2} (0)
$<$ \ce{HClO}, \ce{ClO-} ($+$I) $<$ \ce{ClO3-} ($+$V).
\end{solution}

\begin{solution}{exo:b1:e-ph-diagrams:2}
$-0.059 \times 3/1 = -0.177$; $+0.059 \times 2/2 = +0.059$;
$-0.059 \times 8/5 = -0.094$; $-0.059 \times 2/2 = -0.059$ V per pH unit.
For the copper couple the protons are released on the side of the reduced
form ($m = -2$): the potential rises with pH.
\end{solution}

\begin{solution}{exo:b1:e-ph-diagrams:3}
\ce{H+}/\ce{H2} passes through 0~V at pH 0 and never reaches 0.50~V for
$\mathrm{pH} \geqslant 0$. \ce{O2}/\ce{H2O} reaches 0.50~V at
$\mathrm{pH} = 0.73/0.059 = 12.4$ and 0~V only at pH 20.8, outside the
scale. The band is \qty{1.23}{V} wide at every pH, pH 7 included
(from $-0.41$ to $+0.82$~V).
\end{solution}

\begin{solution}{exo:b1:e-ph-diagrams:4}
(1, 1.0~V): \ce{Fe^3+}. (4, 0): the line \ce{Fe(OH)3}/\ce{Fe^2+} is at
$1.09 - 0.71 = 0.38$~V and the point lies below it, above $-0.47$~V:
\ce{Fe^2+}. (10, $-0.4$~V): between $-0.07 - 0.59 = -0.66$ and $0.28 -
0.59 = -0.31$~V: \ce{Fe(OH)2}. (10, 0.5~V): \ce{Fe(OH)3}.
\end{solution}

\begin{solution}{exo:b1:e-ph-diagrams:5}
\ce{Fe^2+}/\ce{Fe}: $-0.41 + 0.030 \times (-6) = -0.59$~V.
\ce{Fe^3+}/\ce{Fe(OH)3}: $14.00 - \frac13(38.55 - 6) = 3.15$. The domains of
the dissolved ions (corrosion) grow: \ce{Fe^2+} extends down to $-0.59$~V
and \ce{Fe^3+} to pH 3.15; metal and hydroxides lose ground.
\end{solution}

\begin{solution}{exo:b1:e-ph-diagrams:6}
\ce{Fe(OH)3 + 3H+ + e- -> Fe^2+ + 3H2O}: $E = E^\circ + 0.059\log\frac{[\ce{H+}]^3}{[\ce{Fe^2+}]}
= (E^\circ - 0.059\log c) - 0.177\,\mathrm{pH}$. At pH 1.82 the line meets
\ce{Fe^3+}/\ce{Fe^2+} at 0.77~V, so the constant is $0.77 + 0.177 \times 1.82
= 1.09$~V.
\end{solution}

\begin{solution}{exo:b1:e-ph-diagrams:7}
$E^\circ(\ce{Cu+}/\ce{Cu}) = 0.52 > E^\circ(\ce{Cu^2+}/\ce{Cu+}) = 0.16$: by
\cref{prop:b1:e-ph-diagrams:disproportionation-criterion} \ce{Cu+}
disproportionates. Single boundary: $E = 0.34 + 0.030\log 0.010 =
0.28$~V.
\end{solution}

\begin{solution}{exo:b1:e-ph-diagrams:8}
At pH 0 the domain of copper (below 0.28~V) overlaps that of water and
lies above the line of hydrogen: no reaction with hydrochloric acid
without air. Nitric acid: \ce{NO3-}/\ce{NO} at 0.96~V lies above the
domain of copper, which is oxidised:
\ce{3Cu + 2NO3- + 8H+ -> 3Cu^2+ + 2NO + 4H2O}.
% ledger: dfg:NO3-_ao, dfg:NO_g (E0 computed in figdata/bachelor-1/standard-potentials.py)
\end{solution}

\begin{solution}{exo:b1:e-ph-diagrams:9}
At pH 3 the line \ce{O2}/\ce{H2O} is at $1.23 - 0.18 = 1.05$~V, above the
whole domain of \ce{Fe^2+}: oxygen oxidises iron(II). Above pH 1.82 the
iron(III) formed precipitates as the yellow-brown hydroxide:
\ce{4Fe^2+ + O2 + 10H2O -> 4Fe(OH)3 + 8H+}.
\end{solution}

\begin{solution}{exo:b1:e-ph-diagrams:10}
Verticals: $14.00 - \frac12(16.38 - 2) = 6.81$; and, from
\ce{Zn(OH)2 + 2OH- <=> Zn(OH)4^2-} ($K = 10^{-1.93}$) with
$[\ce{Zn(OH)4^2-}] = 10^{-2}$, $[\ce{OH-}]^2 = 10^{-0.07}$, pH 13.97.
\ce{Zn^2+}/\ce{Zn}: $-0.76 + 0.030 \times (-2) = -0.82$~V. \ce{Zn(OH)2}/\ce{Zn}:
slope $-0.059$, through $(6.81, -0.82)$: $E = -0.42 - 0.059\,\mathrm{pH}$.
\ce{Zn(OH)4^2- + 4H+ + 2e- -> Zn + 4H2O}: slope $-0.118$, through
$(13.97, -1.24)$: $E = 0.41 - 0.118\,\mathrm{pH}$.
\end{solution}

\begin{solution}{exo:b1:e-ph-diagrams:11}
The domain of iron lies below the hydrogen line at every pH. pH 2: iron
is oxidised to \ce{Fe^2+} with release of hydrogen (corrosion). pH 7:
the reaction is still possible, but the hydrogen line ($-0.41$~V) is
close to \ce{Fe^2+}/\ce{Fe} and the reaction is very slow. pH 13: iron is
oxidised to \ce{Fe(OH)2}, a \omterm{def:b1:e-ph-diagrams:corrosion-domains}{passivation domain}: a coat forms and the
nail keeps its shine. With dissolved oxygen the line \ce{O2}/\ce{H2O},
far above, oxidises iron to iron(III) at every pH: rust, \ce{Fe(OH)3}
and its dehydrated oxides.
\end{solution}

\begin{solution}{exo:b1:e-ph-diagrams:12}
Zinc has the lower domain: in contact with the steel it is oxidised
first, \ce{2Zn + O2 + 2H2O -> 2Zn(OH)2} at pH 8, and the electrons it
releases reduce the oxygen at the steel surface, which stays metallic.
Magnesium ($E^\circ = -2.36$~V) is lower still and works too; copper
(0.34~V) is above iron: joined to copper, the iron would be the metal
oxidised, faster than alone.
\end{solution}

\begin{solution}{pb:b1:e-ph-diagrams:1}
\textbf{1.} $-$I, 0, $+$I, $+$I.
\textbf{2.} \ce{Cl-} at the bottom, \ce{Cl2}, then \ce{HClO} and \ce{ClO-}
at the top.
\textbf{3.} \ce{HClO} and \ce{ClO-}, an acid--base pair.
\textbf{4.} \ce{Cl2 + 2e- -> 2Cl-}.
\textbf{5.} \ce{2HClO + 2H+ + 2e- -> Cl2 + 2H2O}.
\textbf{6.} \ce{ClO- + H2O + 2e- -> Cl- + 2OH-}.
\textbf{7.} \ce{HClO} below pH 7.55, \ce{ClO-} above.
\textbf{8.} No electron is exchanged; the boundary is fixed by $K_a$ alone.
\textbf{9.} $1/(1 + 10^{7.4 - 7.55}) = 0.59$: \qty{59}{\percent}.
\textbf{10.} $1/(1 + 10^{0.45}) = 0.26$: three quarters of the active
chlorine would be the weaker disinfectant \ce{ClO-}.
\textbf{11.} \ce{ClO-}.
\textbf{12.} $E = 1.396 + 0.0296\log\frac{[\ce{Cl2}]}{[\ce{Cl-}]^2} = 1.396
+ 0.0296\log\frac{0.005}{10^{-4}} = 1.396 + 0.050 = 1.446$~V; no \ce{H+}
in the \omterm{def:b1:nernst:couple}{half-equation}.
\textbf{13.} $E = 1.594 + 0.0296\log\frac{[\ce{HClO}]^2[\ce{H+}]^2}{[\ce{Cl2}]}
= 1.594 + 0.0296\log\frac{10^{-4}}{0.005} - 0.0592\,\mathrm{pH} = 1.594 -
0.050 - 0.0592\,\mathrm{pH}$.
\textbf{14.} Equating, $0.0592\,\mathrm{pH} = 1.594 - 1.396 - 2 \times 0.0503
= 0.0974$, $\mathrm{pH} = 1.646$.
\textbf{15.} \ce{HClO + H+ + 2e- -> Cl- + H2O}: one proton for two
electrons, slope $-0.0592/2 = -0.0296$.
\textbf{16.} Adding the two \omterm{def:b1:nernst:couple}{half-equations} (one electron per chlorine
atom each): $2E^\circ = 1.594 + 1.396$, $E^\circ(\ce{HClO}/\ce{Cl-}) =
1.495$~V; boundary $E = 1.495 - 0.0296\,\mathrm{pH}$.
\textbf{17.} \ce{ClO- + 2H+ + 2e- -> Cl- + H2O}: $E = 1.718 - 0.0592\,\mathrm{pH}$
(constant chosen for continuity: $1.495 - 0.0296 \times 7.55 = 1.272 = 1.718
- 0.0592 \times 7.55$).
\textbf{18.} \ce{Cl-} below the lines; \ce{Cl2} in a small triangle at
$\mathrm{pH} < 1.65$ above 1.446~V; \ce{HClO} above the line of slope
$-0.0296$ up to pH 7.55, \ce{ClO-} beyond. The line \ce{O2}/\ce{H2O},
from 1.23 to 0.40~V, lies below all these boundaries (the computed
diagram of the chapter's script has the same shape).
\textbf{19.} No: its domain lies above the line of oxygen, so it can
oxidise water, but the reaction is slow. A pool loses its chlorine mainly
because it oxidises what swimmers and the air bring in, and through
sunlight; it must be topped up.
\textbf{20.} \ce{Cl2} has no domain there: it disproportionates,
\ce{Cl2 + H2O -> HClO + Cl- + H+}; in base
\ce{Cl2 + 2OH- -> ClO- + Cl- + H2O}.
\textbf{21.} By passing chlorine into sodium hydroxide solution (the
reaction of question 20).
\textbf{22.} \ce{HClO + Cl- + H+ -> Cl2 + H2O}, a \omterm{def:b1:e-ph-diagrams:disproportionation}{comproportionation}.
\textbf{23.} Below pH 1.6, hypochlorite and chloride comproportionate: chlorine forms,
and beyond its \omterm{def:b1:precipitation:solubility}{solubility} escapes as a toxic gas.
\textbf{24.} At $c = \qty{0.010}{mol/L}$, pH 4 is outside the domain of
\ce{Cl2}. But the boundary depends on $c$: repeating questions 12--14 for
a general $c$ gives $\mathrm{pH} = 3.34 + \log(2c)$, which reaches 3.6 for
$c = \qty{1}{mol/L}$, the order of concentration of bleach. Mixing
concentrated products near pH 4 is not safe either.
\textbf{25.} \textbf{pH 1.6}: above it, at $c = \qty{0.010}{mol/L}$,
dissolved chlorine disproportionates into hypochlorous acid and chloride.
\end{solution}
